%document
\documentclass[10pt]{beamer}
%theme
\usetheme{metropolis}
% packages
\usepackage{color}
\usepackage{listings}
\usepackage[ngerman]{babel}
\usepackage[utf8x]{inputenc}
\usepackage{multicol}


% color definitions
\definecolor{mygreen}{rgb}{0,0.6,0}
\definecolor{mygray}{rgb}{0.5,0.5,0.5}
\definecolor{mymauve}{rgb}{0.58,0,0.82}

\lstset{
    backgroundcolor=\color{white},
    % choose the background color;
    % you must add \usepackage{color} or \usepackage{xcolor}
    basicstyle=\footnotesize\ttfamily,
    % the size of the fonts that are used for the code
    breakatwhitespace=false,
    % sets if automatic breaks should only happen at whitespace
    breaklines=true,                 % sets automatic line breaking
    captionpos=b,                    % sets the caption-position to bottom
    commentstyle=\color{mygreen},    % comment style
    % deletekeywords={...},
    % if you want to delete keywords from the given language
    extendedchars=true,
    % lets you use non-ASCII characters;
    % for 8-bits encodings only, does not work with UTF-8
    frame=single,                    % adds a frame around the code
    keepspaces=true,
    % keeps spaces in text,
    % useful for keeping indentation of code
    % (possibly needs columns=flexible)
    keywordstyle=\color{blue},       % keyword style
    % morekeywords={*,...},
    % if you want to add more keywords to the set
    numbers=left,
    % where to put the line-numbers; possible values are (none, left, right)
    numbersep=5pt,
    % how far the line-numbers are from the code
    numberstyle=\tiny\color{mygray},
    % the style that is used for the line-numbers
    rulecolor=\color{black},
    % if not set, the frame-color may be changed on line-breaks
    % within not-black text (e.g. comments (green here))
    stepnumber=1,
    % the step between two line-numbers.
    % If it's 1, each line will be numbered
    stringstyle=\color{mymauve},     % string literal style
    tabsize=4,                       % sets default tabsize to 4 spaces
    % show the filename of files included with \lstinputlisting;
    % also try caption instead of title
    language = Java,
	showspaces = false,
	showtabs = false,
	showstringspaces = false,
	escapechar = ,
}

\def\ContinueLineNumber{\lstset{firstnumber=last}}
\def\StartLineAt#1{\lstset{firstnumber=#1}}
\let\numberLineAt\StartLineAt



\newcommand{\codeline}[1]{
	\alert{\texttt{#1}}
}

% This Document contains the information about this course.

% Authors of the slides
\author{Fritz Windisch, Lukas Tassanyi}

% Name of the Course
\institute{Java-Exercise}

% Fancy Logo
\titlegraphic{\hfill\includegraphics[height=1.25cm]{../templates/fsr_logo_cropped}}


\title{Exercise}
\subtitle{Maven and Audio}
\date{\today}


\begin{document}

\begin{frame}
    \titlepage
\end{frame}

\begin{frame}[fragile]{Maven and Audio}
    Create a small GUI application, which can play sounds/music from a list using Swing, JLayer and Maven.
\end{frame}

\begin{frame}[fragile]{Maven and Audio}
    Use a Maven build environment to fetch the JLayer library. Example pom.xml:
    \begin{lstlisting}
    <?xml version="1.0" encoding="UTF-8"?>
<project xmlns:xsi="http://www.w3.org/2001/XMLSchema-instance"
         xmlns="http://maven.apache.org/POM/4.0.0"
         xsi:schemaLocation="http://maven.apache.org/POM/4.0.0 http://maven.apache.org/xsd/maven-4.0.0.xsd">
    <modelVersion>4.0.0</modelVersion>
    
    <groupId>com.example</groupId>
    <artifactId>soundboard</artifactId>
    <version>1.0-SNAPSHOT</version>
    ...\end{lstlisting}
\end{frame}

\begin{frame}[fragile]{Maven and Audio}
\begin{lstlisting}
    ...
    <dependencies>
        <dependency>
            <groupId>javazoom</groupId>
            <artifactId>jlayer</artifactId>
            <version>1.0.1</version>
        </dependency>
    </dependencies>
</project>\end{lstlisting}
\end{frame}

\begin{frame}[fragile]{Maven and Audio}
    Now you can use JLayer to play .mp3 files:
    \begin{lstlisting}
    Player player = new Player(new FileInputStream(new File("test.mp3"));
    player.play();\end{lstlisting}
    You can also use the pausable player provided in the course materials (you can create a pause button then).
\end{frame}

\begin{frame}{Maven and Audio}
    Tasks:
    \begin{itemize}
        \item Get IntelliJ or Eclipse - BlueJ will not work
        \item Create a new maven project, add JLayer as dependency
        \item Download/Retrieve some of your favorite music/sounds as .mp3
        \item Create a swing gui, which allows the user to play different .mp3 files
        \item (Optional) Allow the user to pause/resume the currently playing file
    \end{itemize}
    You can find GUI element references \href{http://www.java-tutorial.org/swing.html}{here}.
\end{frame}

\end{document}

